\documentclass[
	fontsize=12pt,          % default font size 12pt
	paper=a4,               % DIN A4 page format
	numbers=noenddot,       % remove dots behind chapter numbers (e.g. 1.5 not 1.5.)
	listof=totoc,           % add list of figures, tables, etc. to ToC
	listof=entryprefix,     % add entry name to figures, tables, etc.
	listof=nochaptergap,    % no chapter gap for figures, tables, etc.
	bibliography=totoc,     % add bibliography to ToC but without a chapter number
	parskip=half            % half line spacing between paragraphs
	openany                 % chapters can start on any page
]{scrbook}
\usepackage{tcolorbox}

\include{../../for_latex/preamble}
\title{\Huge\textbf{Stabkräfte im Fachwerk}}
\author{Luis Staudt}
\date{}

% Code-Highlighting für Java
\definecolor{codegray}{rgb}{0.5,0.5,0.5}
\definecolor{codepurple}{rgb}{0.58,0,0.82}
\definecolor{backcolour}{rgb}{0.95,0.95,0.92}
\definecolor{keywordcolor}{rgb}{0.8,0.47,0.13}

\lstdefinestyle{java}{
	language=,
	basicstyle=\ttfamily\small,
	keywordstyle=\bfseries\color{blue!70!black},
	stringstyle=\color{red!60!black},
	commentstyle=\itshape\color{gray},
	numbers=left,
	numberstyle=\tiny\color{gray},
	numbersep=8pt,
	frame=single,
	framerule=0.4pt,
	rulecolor=\color{gray!50},
	breaklines=true,
	showstringspaces=false,
	tabsize=4,
	xleftmargin=1.5em,
	framexleftmargin=1.5em,
	aboveskip=0.8em,
	belowskip=0.6em
}
\lstset{style=java}

\newtcolorbox{hinweis}[1][Hinweis]{
	colback=blue!5,
	colframe=blue!40!black,
	fonttitle=\bfseries,
	title=#1,
	boxrule=0.5pt,
	arc=2pt
}

\newtcolorbox{beispiel}[1][Beispiel]{
	colback=green!5,
	colframe=green!40!black,
	fonttitle=\bfseries,
	title=#1,
	boxrule=0.5pt,
	arc=2pt
}

\begin{document}

% --- Titelblock ---
\begin{center}
	\rule{\textwidth}{0.8pt}\\[0.6em]
	{\LARGE\bfseries Intensivtutorium}\\[0.3em]
	{\large Programmierung}\\[0.6em]
	\begin{tabular}{r l @{\hspace{2cm}} r l}
		\textbf{Semester:}      & Sommersemester 2026                    & \textbf{Tutor:}         & Luis Staudt \\
		\textbf{Studiengänge:}  & MIB, OMB, PEB, MVB         & & \\
	\end{tabular}\\[0.6em]
	\rule{\textwidth}{0.8pt}
\end{center}
\vspace{1em}

% --- Seitenkopf und -fuß (ab Seite 2) ---
\pagestyle{fancy}
\fancyhf{}
\lhead{\small Intensivtutorium -- Grundlagen der Programmierung}
\rhead{\small SS 2026}
\cfoot{\thepage}
\renewcommand{\headrulewidth}{0.4pt}


% =============================================
% Übungsaufgabe -- Stabkräfte im Fachwerk
% =============================================

\section*{Übungsaufgabe -- Stabkräfte im Fachwerk}

Gegeben ist die aus Vorlesung und Übung bekannte Datei \texttt{fachwerk.jar}.
Nach der Berechnung sollen für ein gewähltes Fachwerk die \textbf{Stabkräfte} ausgewertet werden.

Ein Stab steht unter \textbf{Zug}, wenn seine Stabkraft $> 0$ ist, und unter \textbf{Druck}, wenn sie $< 0$ ist.
Ist die Stabkraft (im Rahmen einer Epsilon-Umgebung) praktisch null, handelt es sich um einen \textbf{Nullstab}.

\begin{center}
\begin{tikzpicture}[>=Stealth, scale=1.0,
	knoten/.style={circle, fill=black, inner sep=1.6pt},
	stab/.style={thick, gray!70!black}]
	\coordinate (A) at (0,0); \coordinate (B) at (2.5,0); \coordinate (C) at (5,0);
	\coordinate (D) at (1.25,2.0); \coordinate (E) at (3.75,2.0);
	\draw[stab] (A)--(B); \draw[stab] (B)--(C);
	\draw[stab] (A)--(D); \draw[stab] (D)--(B); \draw[stab] (D)--(E);
	\draw[stab] (E)--(B); \draw[stab] (E)--(C);
	\foreach \p in {A,B,C,D,E} \node[knoten] at (\p) {};
	\draw[->, thick] (B)--++(0,-0.9);
	\node[below, font=\footnotesize] at ($(B)+(0,-0.9)$) {Last};
\end{tikzpicture}\\[-0.3em]
{\footnotesize\itshape Schematische Darstellung eines Fachwerks.}
\end{center}

\begin{hinweis}[Physikalischer Hintergrund]
	Ein Zugstab wird auseinandergezogen (positive Stabkraft), ein Druckstab wird gestaucht (negative Stabkraft).
	Nullstäbe tragen im belasteten Zustand keine Kraft; sie sind konstruktiv aber wichtig, weil sie das Fachwerk aussteifen und ein Ausknicken verhindern.
	Der betragsmäßig am stärksten beanspruchte Stab bestimmt die Dimensionierung des Bauteils.
\end{hinweis}

\subsection*{Aufgabe}

Ergänzen Sie das Programm so, dass nach der Berechnung folgende Auswertungen erfolgen:

\begin{enumerate}[label=\textbf{\arabic*.)}]
	\item Zählen Sie, wie viele Stäbe \textbf{Zugstäbe}, \textbf{Druckstäbe} bzw.\ \textbf{Nullstäbe} sind, und geben Sie die drei Anzahlen aus.

	\item Bestimmen Sie den \textbf{am stärksten beanspruchten Stab}, also den Stab mit dem betragsmäßig größten Stabkraftwert.
	Geben Sie dessen Nummer, seine beiden Knotennummern sowie den Stabkraftwert aus.
\end{enumerate}

\begin{hinweis}[Epsilon-Umgebung beachten]
	Wegen Gleitkomma-Rundung ist \glqq null\grqq{} über eine Epsilon-Umgebung zu prüfen: $|F| < \varepsilon$.
	Benötigte Methoden: \texttt{gibAnzahlStaebe()}, \texttt{gibStab(int i)}, \texttt{gibStabkraft()},\\ \texttt{gibKnotenNrA()}, \texttt{gibKnotenNrE()}.
	Der Betrag lässt sich mit \texttt{Math.abs(\dots)} bilden.
\end{hinweis}

\medskip
Vervollständige das folgende Programmgerüst:

\begin{lstlisting}
public class TrussBarForces {
    public static void main(String[] args) {
        System.out.println("Welches Fachwerk wollen Sie rechnen (0-4)?");
        int wahl = Tastatureingabe.leseInt();
        Fachwerk fw = new Fachwerk(wahl);
        fw.berechneGroessen();
        double epsilon = 1e-6;
        // hier erweitern: Zug-/Druck-/Nullstaebe zaehlen sowie
        // den am staerksten beanspruchten Stab bestimmen
    }
}
\end{lstlisting}

\begin{beispiel}[Mögliche Ausgabe (Fachwerk~2)]
\begin{lstlisting}[language={}]
Zugstaebe:   18
Druckstaebe: 17
Nullstaebe:  2
Am staerksten beansprucht: Stab 5 (Knoten 5 -> 6) mit -120.5985074534
\end{lstlisting}
\end{beispiel}

\end{document}
